\chapter{Angewandte Sanitätshilfe ---
\gAasRsN{RS~XIII}}

\author{Gabriel, Hochreiter, Motal}

\minitoc

\section{Lagerungsarten}\label{topic:lagerungsarten}\index{Lagerungsarten}

%========Beginn Abschnitt Emhofer============================================

\subsection{{\quotedblbase}Hilfe, ein Patient!{\textquotedblleft} - Was
mach{\textquotesingle} ich jetzt?}

Gerade zu Beginn der T\"atigkeit als Rettungssanit\"ater f\"uhlt man sich beim Eintreffen am Notfallort durch die große Menge an Sinneseindr\"ucken und Informationen bald \"uberlastet. Um Sie sicher durch diese Informationsflut zu leiten, wollen wir Ihnen in diesem Abschnitt ein paar Anhaltspunkte pr\"asentieren, welche den Einsatzablauf strukturieren und Ihnen und somit auch dem Patient Sicherheit geben.
Vorgehen am Notfallort:
\begin{itemize}
\item Bevor es los geht: Selbstschutzmaßnahmen
\item 1. Schritt: \gEs{Ersteinsch\"atzung des Patienten}
\item 2. Schritt: \gEs{Erstmaßnahmen}
\item 3. Schritt: \gEs{Durchatmen, nachdenken, planen}
\end{itemize}
W\"ahrend die ersten zwei Schritte ein gezieltes Vorgehen erfordern, bietet der dritte Schritt Spielraum für die individuelle Untersuchung und Befragung des Patienten entsprechend den Symptomen. Im Folgenden stellen wir Ihnen die Schritte im Einzelnen vor.

\subsection{Reden - Die Anamnese erz\"ahlt uns die
Geschichte des Patienten}

Neben der k\"orperlichen Untersuchung gibt das Gespr\"ach mit dem Patienten wertvolle Hinweise auf das Notfallgeschehen. Diese Erhebung der Kranken- bzw. Unfallgeschichte nennt man \gEs{Anamnese}. Eine gute Anamnese ist mindestens genauso wichtig wie die Erhebung der Befunde! Leider wird in der Praxis oft der Messwerterhebung eine gr\"oßere Bedeutung beigemessen.

Auch das Anamnesegespräch soll strukturiert durchgef\"uhrt werden. Daf\"ur stehen mehrere Abfrageschemata zur Verf\"ugungen, von welchen wir Ihnen ein in \"Osterreich weit verbreitetes vorstellen m\"ochten: Das \gEs{S-KAMEL-Schema}. Nach einer Einstiegsfrage, z.B. "`Was können wir für Sie tun?"', beinhaltet das S-KAMEL-Schema folgende Punkte:
\begin{itemize}
\item Symptome und Schmerzen
\item Krankheiten 
\item Allergien 
\item Medikamente 
\item Ereignisse vor Notfalleintritt 
\item Letzte...
\end{itemize}
Dieses Abfrageschema kann Ihnen im hektischen Notfall Sicherheit geben. Durch das Abarbeiten der Fragen haben Sie die Sicherheit, dass keine wichtigen Informationen vergessen werden. Außerdem ist es für den Patienten vertrauensbildend, da die gezielte Befragung professionell wirkt. Im Folgenden m\"ussen wir uns die einzelnen Punkte des Schemas genauer ansehen.

\gFazit{S{\textquotesingle}KAMEL denkt an alles}

\subsubsection{Symptome und Schmerzen}
Die Einstiegsfrage hat uns u.U. bereits auf das Hauptsymptom hingewiesen. Ist dies nicht der Fall, so muss als erstes nach dem Hauptsymptom gefragt werden. Fragen können z.B. sein:
\begin{itemize}
\item Welche Beschwerden haben Sie?
\item Wo drückt der Schuh?
\end{itemize}
Außerdem muss immer auch speziell nach Schmerzen gefragt werden!
\begin{itemize}
\item Haben Sie Schmerzen?
\end{itemize}
Gibt der Patient an Schmerzen zu haben, m\"ussen wir weitere Informationen erfragen, damit wir uns ein klares Bild \"uber die Beschwerden machen können:
\begin{itemize}
\item \gEs{Lokalisation}: Wo sind die Schmerzen?
\item \gEs{Zeitdauer/Beginn}: Seit wann haben Sie die Schmerzen? Wie lange haben Sie die Schmerzen schon? Haben Sie so einen Schmerz schon einmal gehabt? Haben die Schmerzen plötzlich oder langsam eingesetzt?
\item \gEs{Qualität}: Wie ist der Schmerz? Bohrend, stechend, ziehend, krampfartig?
\item \gEs{Ausstrahlung}: Strahlt der Schmerz irgendwohin aus?
\item \gEs{Provokation}: Was macht den Schmerz besser oder schlechter?
\item \gEs{Stärke}: Wie stark ist der Schmerz? (In einigen Lehrbüchern wird empfohlen, den Schmerz anhand einer Skala von 1-10 beurteilen zu lassen, wir finden jedoch, dass auch eine solche Beurteilung recht subjektiv ist und daher nur mit Vorsicht verwendet werden sollte. Setzen Sie die Aussage des Patienten immer in Relation zu seinem Gesichtsausdruck und zu seiner Körperhaltung! Einem Patient mit starken Schmerzen sieht man dies an!)
\end{itemize}

\subsubsection{Krankheiten}
Notfälle ereignen sich manchmal aufgrund der akuten Verschlechterung einer bestehenden Vorerkrankung. Außerdem hilft die Beurteilung der Messwerte, wenn man über die Grunderkrankungen des Patienten Bescheid weiß. Zum Beispiel kann die Ursache eines hohen Blutdrucks eine bereits bestehende Erkrankung aber auch ein Notfallgeschehen sein. Neben der Befragung kann auch ein alter Spitalsbericht (Arztbrief) Auskunft über bestehende Vorerkrankungen liefern. Allgemeine Fragen, wie z.B.
\begin{itemize}
\item Leiden Sie unter irgendwelchen Krankheiten?
\item Bestehen bei Ihnen Vorerkrankungen?
\end{itemize}
können anschließend durch symptom-spezifische Fragen ergänzt werden:
\begin{itemize}
\item Hatten Sie schon einmal einen Herzinfarkt?
\item Leider Sie unter hohem Blutdruck?
\item Sind Sie zuckerkrank?
\end{itemize}

\subsubsection{Allergien}
Auf die Frage nach möglichen Allergien wird häufig vergessen, ein Informationsmangel kann aber fatale Auswirkungen haben. Fragen Sie daher immer nach Allergien, in einer zweiten Frage sogar speziell nach Medikamentenallergien! Damit können Sie für Not- oder Spitalsarzt wichtige Informationen erheben. (Natürlich könnten das die Ärzte auch selbst fragen, aber denken Sie daran, dass ein Notfallpatient bewusstlos werden kann und es deshalb wichtig ist, dass sie diese Informationen einholen, solange der Patient noch ansprechbar ist!)

\subsubsection{Medikamente}
Die vom Patienten eingenommenen Medikamente lassen Rückschlüsse auf seine Grunderkrankungen zu und sind im Falle eines längeren Spitalsaufenthalts eine wertvolle Information. Auch die Regelmäßigkeit oder Unregelmäßigkeit der Einnahme kann Ihnen Anhaltspunkte für das Finden Ihrer Verdachtsdiagnose bzw. der Umstände des Patienten liefern. Fragen können sein:
\begin{itemize}
\item Welche Medikamente nehmen Sie?
\item Wogegen nehmen Sie diese Medikamente?
\item Nehmen Sie Ihre Medikamente regelmäßig?
\item Wann haben Sie Ihre Medikamente zuletzt genommen?
\end{itemize}

\subsubsection{Ereignisse vor Notfalleintritt}
Oft ereignen sich vor einem Notfall Schlüsselereignisse, welche den Notfall herbei führen oder diesen auslösen. Fragen Sie daher immer nach möglichen Besonderheiten oder Tätigkeiten des Patienten vor dem Notfall:
\begin{itemize}
\item Was haben Sie vor Beginn der Symptome gemacht?
\item Haben Sie sich vor Beginn der Symptome besonders angestrengt?
\item u.v.m.
\end{itemize}
In dieser Kategorie kann es auch hilfreich sein, wenn Sie Angehörige oder Zeugen befragen:
\begin{itemize}
\item Hat der Patient vor dem Notfall etwas Eingenommen (Überdosis)?
\item Hat der Patient gekrampft bevor er bewusstlos liegen geblieben ist?
\item Ist der Patient vor dem Eintreten der Symptome gestürzt?
\item u.v.m.
\end{itemize}

\subsubsection{Letzte ...}
Der letzte Punkt des Abfrageschemas bezieht sich auf letzte Ereignisse, welche mit dem Notfall in Zusammenhang stehen könnten. Damit Sie hier die richtigen Fragen stellen können, ist ein fundiertes Fachwissen notwendig. Abhängig von den Beschwerden können z.B. folgende Fragen zielführend sein:
\begin{itemize}
\item Wann haben Sie zuletzt etwas gegessen/getrunken? (Ist der Patient nüchtern?)
\item Wann hatten Sie zuletzt Stuhlgang? (weiterführende Frage: Wie hat der Stuhl ausgesehen?)
\item Wann war die letzte Regelblutung? (Besteht die Möglichkeit einer Schwangerschaft?)
\item Wann waren Sie zuletzt beim Arzt bzw. im Spital?
\end{itemize}

\subsection{Kommunikationsregeln}
Der Erfolg Ihres Anamnesegespräches ist für das weitere Vorgehen u.U. entscheidend. Oft erhalten Sie wichtige Information zum Zustand des Patienten nicht durch die gemessenen Werte sondern durch die Befragung! Der Erfolg des Gespräches hängt dabei einerseits von ihrem kommunikativem Geschick, andererseits aber auch von der Bereitschaft des Patienten ab. Im Folgenden möchten wir Ihnen ein paar wertvolle Tipps geben, wie Sie die Kommunikation so gestalten können, dass Sie die gewünschten Informationen erhalten:
\begin{itemize}
\item Stellen Sie sich mit Ihrem Namen vor und - so es der Notfall zulässt - geben Sie dem Patienten die Hand!
\item Begeben Sie sich auf die gleiche Ebene (auf Augenhöhe) mit dem Patienten!
\item Hören Sie dem Patienten aktiv zu! Geben Sie dem Patienten das Gefühl, dass er verstanden wird, indem Sie z.B. zustimmend Nicken oder die wichtigen Informationen wiederholen.
\item Verwenden Sie eine leicht verständliche Sprache und vermeiden Sie Fachausdrücke!
\item Informieren Sie den Patient über Ihr Vorgehen!
\item Versuchen Sie die Fragen des Patienten mit bestem Wissen und Gewissen zu beantworten! Lügen Sie den Patienten nicht an!
\item Im Sanitäterteam sollten nicht alle Sanitäter auf den Patient einreden! Der Patient sollte nur einen Ansprechpartner haben!
\end{itemize}

Bei Notfällen mit Kindern kann die Kommunikation besonders herausfordernd werden, da Kinder die Angst haben sich oft gegenüber Fremden verschließen. Die oben genannten Regeln gelten daher bei Kindern umso mehr. Zusätzlich können folgende Tipps hilfreich sein:
\begin{itemize}
\item Kinder nehmen non-verbale Signale stärker wahr als Erwachsene. Versuchen Sie daher durch einen freundlichen Gesichtsausdruck und eine offene Körperhaltung Vertrauen zu bilden!
\item Erklären Sie dem Kind (altersgerecht) warum Sie gerufen wurden!
\item Sprechen Sie das Kind direkt und mit seinem Namen an! Vermeiden Sie es über das Kind in der 3. Person zu sprechen, damit das Kind nicht das Gefühl bekommt, dass hinter seinem Rücken getuschelt wird. Dies würde dem Kind unnötig Angst machen!
\item Seien Sie immer ehrlich! Kinder merken sehr schnell, wenn Sie lügen. Lügen Sie ein Kind an, so wird das Kind schnell das Vertrauen in Sie verlieren und jede weitere Untersuchung verweigern!
\item Geben Sie dem Kind ein Spielzeug in die Hand!
\item Trennen Sie das Kind niemals von der Mutter!
\item Hat das Kind Angst vor einer bestimmten Untersuchung oder einer Maßnahme, zeigen Sie diese zuerst an der Mutter oder an einem Stofftier vor!
\item Loben Sie das Kind nach der Untersuchung bzw. nach der Maßnahme!
\item Da Kinder ihre Beschwerden nicht eindeutig beschreiben können, lassen Sie sich z.B. den Ort der Schmerzen zeigen.
\end{itemize}

%========Ende Abschnitt Emhofer============================================


\section{Untersuchung und Anamnese}

Hier lernen Sie wie Sie die bereits gelernten therotischen Grundlagen in
der Praxis einsetzen, wie und in welcher Reihenfolge Sie im Einsatz
vorgehen, wie Sie Dringlichkeiten einsch\"atzen und Verdachtsdiagnosen
erarbeiten k\"onnen und wie Sie dabei dann auch noch den \"Uberblick
behalten k\"onnen.

\subsection{{\quotedblbase}Hilfe, ein Patient!{\textquotedblleft} - Was
mach{\textquotesingle} ich jetzt?}
Blabla

\"Ubersicht 

Ersteinsch\"atzung ${\rightarrow}$  Erstma{\ss}nahmen ${\rightarrow}$ 
Plan zurechtlegen wie es weiter gehen soll

blabla

mehr blabla

\subsection{Reden - Die Anamnese erz\"ahlt uns die
Geschichte des Patienten}

\gFazit{S{\textquotesingle}KAMEL denkt an alles}

\subsubsection{Symptome und Schmerzen}
\subsubsection{Krankheiten}



\subsubsection{Allergien}
\subsubsection{Medikamente}
\subsubsection{Ereignisse}
\subsubsection{Das Letzte ...}

\subsection{Untersuchungen k\"onnen uns weiterhelfen}

\subsubsection{Sehen -- H\"oren -- F\"uhlen }
Die Seh-, Hör- und Tastsinne sind bei der körperlichen
Untersuchung unverzichtbar. 

Zuerst wird der Patient und sein
Körper betrachtet. Wir versuchen uns bereits bekannte Symptome wahrzunehmen. Hierbei können uns viele
augenscheinliche Sachen auffallen: zum Beispiel
\gE{Fettleibigkeit}, \gE{Zyanose}, \gE{Blässe},
\gE{Schweiß}, \gE{Fremdkörper}, \gE{Verletzungen}, \gE{Fehlstellungen},
\ldots.

Der nächste Schritt betrifft das Hören: Hier sind zum
Beispiel die Atemgeräusche einzuordnen, die man eventuell schon ohne
Hilfsmittel vernehmen kann (brodelndes Atemgeräusch beim
Lungenödem, pfeifendes Atemgeräusch beim Asthmaanfall,
röcheln bei Atemwegsverlegung, \ldots)

Wenn wir gesehen und gehört haben können wir versuchen zu
fühlen: Wir können den (radialis-) Puls fühlen, können
(stehende) Hautfalten beurteilen oder das Abdomen abtasten.

\paragraph{Abtasten des
Abdomens}\label{topic:untersuchung-abdomen-abtasten}

Das Abdomen kann man abtasten um die Anspannung der
Bauchdecke zu beurteilen (sie voraus), Schmerzen zu
beurteilen und zu zuordnen, Fremdkörper zu ertasten und
anderes. Man beginnt damit den Patienten zu fragen ob er ihm
auch Schmerzen hat und wo er diese hat. Die Untersuchung
beginnt man in dem Quadranten der am weitesten von den
Schmerzen entfernt ist!

Man legt nun eine Hand flach auf den Bauch des Patienten und
die andere Hand flach auf die erste. Mit den Fingern drückt
man nun vorsichtig mehrere Zentimeter auf den Bauch und
beurteilt auf der Bauch leicht ist oder der Patient
einsparen. Dabei fragt man den Patienten ob dies weh tut.
Dies wiederholt man in allen vier Quadranten.

Weiters kann man beurteilen ob der Patient Schmerzen durch
den Druck hat (Druck Schmerz) oder ob er Schmerzen beim
loslassen bekommt (los lass Schmerz).
\subsubsection{Vitalwerte}
Das erheben von vital werten ist eine Routinetätigkeit und
sollte bei jedem Notfallpatienten erfolgen. Es gibt nur wenig
Situationen bei denen man darauf verzichten kann (zum
Beispiel psychiatrische Patienten).

Die Messung der Herzfrequenz erfolgt mittels Puls fühlen oder
mittels Pulsoxymeter. Beim Puls fühlen wird für 15~Sekunden
der Puls gefühlt und die Pulsschläge gezählt. Anschließend wird
das Ergebnis mit vier multipliziert, man erhält die
Pulsschläge pro Minute. Bei der Pulsoxymetrie wird der
entsprechende Wert abgelesen. Es ist zu beachten was bei
Notfallpatienten die Pulsoxymetrie oft nicht funktioniert (ZB
Zentralisationbei Schock, kalte Finger,\ldots). Neben der
Frequenz wird auch die Regelmäßigkeit beurteilt, das heißt ob
der Puls rhythmisch oder arrhythmisch ist.

Die Messung des Blutdrucks ist im Kapitel Sanitäts Technik
näher beschrieben (s. Kap.\ref{topic:rr-messung}, S.\pageref{topic:rr-messung}.

\subsubsection{Atmung}



\subsubsection{Pulsoxymetrie}

Die Pulsoxymetrie ist eine einfache apparative
Untersuchungsmethode, welche mithilfe eines Pulsoxymeters
durchgeführt wird. Dabei wird ein Sensor auf den Finger des
Patienten geklipt, welcher Licht ausstrahlt. Dieses Licht
wird von einem anderen Sensor aufgefangen, das Gerät kann
aus den daraus gewonnenen Informationen auf den Sauerstoffgehalt des
Blutes und die Pulsfrequenz schließen. Dies setzt jedoch
voraus das der Patient in der Peripherie über eine
ausreichende Durchblutung verfügt, da das Gerät sonst keine
Informationen gewinnen kann (mangels Blut zur
Beurteilung). Weitere Hindernisse sind z.B. kalte
Finger (verringerte Durchblutung) oder Nagellack (Licht kann
den Nagellack nicht durchdringen, eventuell mit
Nagellack-Entferner reinigen).

Wie Sauerstoffsättigung wird in Prozent angegeben und wird
mit $SpO_2$ abgekürzt. Der Normalwert beträgt 95-100 %. Bei
Patienten mit COPD findet man oft sehr niedrige Werte (ZB
90\,%).

Bei der Interpretation des $SpO_2$ muss man sehr
vorsichtig sein! Auch ein Normalwert bedeutet nicht dass der Patient
genug Stoitschkow hat beziehungsweise dass keine sorgt
Aufgabe erfolgen soll. Die sorgt Aufgabe richtet sich immer
nach der Diagnose und dem klinischen Zustandsbild des
Patienten, niemals!! Nach dem SP $O_2$-Wert. Der
Sp$O_2$-Wert soll immer im Verlauf beurteilt werden.

Eine besondere Fehlerquelle ist die Kohlenmonoxid-
($CO$)-Vergiftung, bei ihr wird nicht nur der Patient
Schweinchen rosa obwohl er eigentlich zu wenig Sauerstoff
hat, auch das Pulsoxymeter verwechselt Kohlenmonoxid mit
Sauerstoff und zeigt falsch hohe Werte an!

Das Pulsoxymeter eignet sich sehr gut für das Monitoring,
ersetzt aber das wachsame Auge auf den Patienten niemals!

\subsubsection{Blutzuckermessung}

\subsubsection{K\"orpertemperatur}

\subsubsection{Neurocheck}

\subsubsection{Traumacheck}
Suche nach Verletzungen oder Verletzungsanzeichen

\begin{itemize}
\item Patient orientiert? Unfallhergang erinnerlich? Schmerzen?

\end{itemize}
\begin{itemize}
\item Abtasten mit festem Griff: von Kopf bis Fu{\ss} (Sch\"adeldach,
Hinterkopf, Stirn, Wangenknochen, Nasenbein, Kiefer, Hals, Schultern,
Arme, H\"ande, Brustkorb vorne und seitlich (Prellmarken?
Atemmechanik?), Bauch in 4 Quadranten (Prellmarken?
Abwehrspannung/harte Bauchdecke?), Schambein, Becken (stabil?), Beine
(sichtbare Verletzungen?), F\"u{\ss}e: Schuhe und Socken ausziehen)

\item {\quotedblbase}DMS{\textquotedblleft}: Finger und Zehen: Kontrolle
von Motorik (bewegen lassen), Sensibilit\"at ({\quotedblbase}Sp\"uren
Sie das?{\textquotedblleft}) und Durchblutung:
Nagelbettprobe/Rekapillarisierungszeit: leichter Druck auf Nagelende
bis Nagelbett wei{\ss} ist, loslassen. Zeit bis Nagelbett wieder
rosarot ist sollte nicht l\"anger als 1-2 Sekunden sein.
Seitenvergleich! Eine verl\"angerte Rekapzeit deutet auf eine
Durchblutungsst\"orung hin.

\item Kontrolle von Mund, Nase und Ohren auf Blutungen (mit einer
Lampe). Tupfertest: etwas Blut aus Ohren mit einem Tupfer auffangen.
Sollte Liquor im Blut enthalten sein bildet sich um den roten Blutfleck
ein heller, bernsteinfarbener Hof.

\end{itemize}
bei Sch\"adelverletzungen ist ein Neurocheck auf jeden Fall
erforderlich!

Tarumavortr\"age m\"ussen hierher verweisen


\subparagraph{Abtasten des Abomens }
\clearpage
\subsection{Vorgehen in der Praxis}

\paragraph{Szene\"uberblick}
\paragraph[Die ersten \ \ 30 Sekunden ]{Die ersten  30 Sekunden }
Ersteinsch\"atzung ${\rightarrow}$  Erstma{\ss}nahmen ${\rightarrow}$ 
Plan zurechtlegen wie es weiter gehen soll

\subsection{Erster Schritt: Ersteinsch\"atzung}
\begin{center}
\begin{tabulary}{\linewidth}[H]{CL}
\toprule\midrule
 \gEs{Ersteindruck}
 &
Aussehen, Gesichtsausdruck, Haltung, Sprache

Bl\"asse, Zyanose, Schwei{\ss}, Schohnhaltung, Atemnot, verwaschene
Sprache, Des\-orientiertheit, Kr\"ampfe 
\\\hline
 \gEs{Bewusstsein}
 &
Reaktion auf verbale Ansprache, Weckversuch, Schmerzreiz

klar, somnolent, sopor\"os, komat\"os/bewusstlos
\\\hline
 \gEs{Atemweg}
 &
Atemger\"ausch, Gestik

frei oder verlegt
\\\hline
 \gEs{Atmung}
 &
Frequenz, Atemmuster, Anstrengung beim Atmen, Atemger\"ausche,
Heimsauerstoff
\\\hline
 \gEs{Kreislauf}
 &
Bewusstsein, Hautfarbe, Radialispuls, Rekap-Zeit
\\\hline
 \gEs{Traumacheck 1}
 &
Schwere Verletzungen: Backen-, Oberschenkel-, Serienrippenfraktur,
massive Blutung
\\\hline
 \gEs{Hauptbeschwerde}
 &
Grund der Berufung, Schmerzen und andere f\"ur den Patienten dringliche
Symptome
\\\hline
\end{tabulary}
\end{center}
\marginpar{{\sffamily\color{red}\bfseries Alarmzeichen}

\gEs{Massive St\"orungen von} 

Bewusstsein

Atemweg / Atmung 

Kreislauf 

\gEs{Typische Symptome} 

Atemnot die sich nicht bessert 

Brodelndes Atemger\"ausch 

Thoraxschmerz  

Schocksymptome 

entgleiste Vitalwerte

schwere Verletzungen

{\dots} 
}
Nun stellen sich folgende Fragen:

\begin{enumerate}
    \item Was ist das Problem des Patienten?
    \item Ist der Patient akut vital bedroht?
    \begin{itemize}
        \item Anzeichen einer vitalen Bedrohung k\"onnen u.a. sein: St\"orungen
        von Be\-wusstsein, Atemweg, Atmung oder Kreislauf sowie typische
        Symptome wie sich nicht bessernde Atemnot, brodelndes
        Atem\-ge\-r\"ausch, Thoraxschmerz, Schocksymptome und entgleiste
        Vital\-werte.
    \end{itemize}
\end{enumerate}



W\"ahrend des weiteren Verlaufs (Anamnesegespr\"ach, etc.) kann man auf
weitere Befunde sto{\ss}en welche auf eine vitale Be\-drohung hinweisen
(zB. massiv erh\"ohter Blutdruck bei bei vor\-ge\-sch\"adigtem Herzen,
pl\"otzliche oder schleichende Zustands\-verschlechterung, etc.). 

Die Ersteinsch\"atzung muss st\"andig \"uberpr\"uft werden. Neue Befunde
und der Zustand des Patienten k\"onnen die Situation \"andern.

\subsection{Zweiter Schritt: Erstma{\ss}nahmen}

Welche Ma{\ss}nahmen sind f\"ur den Patienten am wichtigsten und
m\"ussen zuerst erfolgen? 

\marginpar{\gH{Bei der Pr\"ufung "`blind"' auf
Nummer sicher gehen und alle gleich f\"ur "`vital bedroht"'
erkl\"aren gilt nicht!}

Viele Patienten bei der Pr\"ufung sind nicht vital bedroht und brauchen
keinen Notarzt. 

Erkl\"art man einen Patienten f\"ur vital bedroht muss man das
begr\"unden k\"onnen!
}


Bei \gEs{vital bedrohten Patienten} gibt es \gEs{5
Standardma{\ss}nahmen} die sofort zu erfolgen haben. Bei
allen anderen Patienten muss individuell \"uberlegt werden welche Ma{\ss}nahmen zuerst zu treffen sind.

\begin{center}
\begin{tabulary}{\linewidth}{LL}
\toprule
\multicolumn{2}{c}{\gEs{Erstma{\ss}nahmen bei vital
bedrohten Patienten}}
\\\midrule
Situationsgerechte Lagerung
 &
\\\midrule
Sauerstoffgabe
 &
je nach Indikation, Vorsicht bei COPD und Hyper\-ventilation
\\\midrule
Notarzt{}-Nachforderung
 &
mit kurzer Begr\"undung
\\\midrule
Reanimationsbereitschaft
 &
Platz schaffen, Beatmungsbeutel zu\-sammen\-bauen, ggfs. Defi und
Absaugung holen und in Griffweite stellen
\\\midrule
Vitalwerte erheben und Monitoring
 &
RR, HF, Pulsoyxymetrie wenn vor\-handen
\\\bottomrule
\end{tabulary}
\end{center}

\subsection{Dritter Schritt: Durchatmen, nachdenken, planen}

Nach den Erstma{\ss}nahmen muss man nun planen welche Ma{\ss}nahmen,
Untersuchungen und Fragen am wichtigsten sind und zuerst durchgef\"uhrt
werden m\"ussen und was eher in der Abfolge nach hinten ver\-schoben
wird. Die folgende Checkliste in Tabelle
\ref{tab:checkliste-patientenmanagement} gibt einen
\"Uberblick \"uber die Ma{\ss}nahmen und das Anamnese\-ge\-spr\"ach, die Reihenfolge ist bei jedem Patienten anders!

%%%%%%%%%%%%%%%%%%%%%%%%%%%%%%%%%%%%%%%%%%%%%%%%%%%%%%%%%%%
\begin{table}
\captionabove{Checkliste Patientenmanagement -- M\"ogliche Ma{\ss}nahmen, die Reihenfolge
ist je nach Patient unterschiedlich}
\label{tab:checkliste-patientenmanagement}

\begin{tabulary}{\linewidth}[H]{CCL}
\toprule
SELBSTSCHUTZ
%\includegraphics[width=1.005cm,height=0.9cm]{AASdev20090228-img14.gif}
 &
 &
Schutzausr\"ustung (Handschuhe, Uniform, Helm)
Gefahrenbereich, Unfallstelle absichern, Hunde wegsperren lassen
ggfs. R\"uckzug, Spezialkr\"afte anfordern, {\dots}
{\dots}
\\\midrule
ERKENNEN

\includegraphics[width=1cm]{icons/eye.png}
 &
 &
Situation/Milieu {\tiny(zB Alter, soziales Umfeld, SMV,
Alkohol, Pensionistenheim, ...)}

offensichtliche Symptome {\tiny(zB Zyanose,
Bewusstlosigkeit, starke Blutung, ...)}

Schockgefahr {\tiny(Zyanose, kalter Schwei{\ss},
Radialispuls tastbar/nicht tastbar, Puls beschleunigt)}
\\\midrule
ERSTE HILFE

\includegraphics[width=1cm]{icons/firstAid.png}
 &
 &
Lagerung

Helmabnahme

Blutstillung

psychischer Beistand

W\"armeerhaltung
\\\midrule
{\multirow{6}{2cm}{\includegraphics[width=1cm]{icons/exclamationMark.png}
ANAMNESE}}
 &
{\sffamily\textbf S}
 &
Symptome \& Schmerzen 

{\tiny(Zeit/Verlauf, St\"arke, Art/Qualit\"at,
Lokalisation/Ausstrahlung, Provokation)}
\\
 &
{\sffamily\textbf K}
 &
Krankheiten {\tiny(Bestehen Grunderkrankungen, zB
Diabetes?)}
\\
 &
{\sffamily\textbf A}
 &
Allergien/Allergische Reaktion (auch
Allergien auf Medikamente)
\\
 &
 {\sffamily\textbf M}
 &
Medikamente {\tiny(Nimmt Patient Medikamente?
Wogegen sind diese?)}
\\
 &
 {\sffamily\textbf E}
 &
Ereignisse vor Notfalleintritt {\tiny(zB
Anstrengung vor Brustschmerz, Sturz, ...)}
\\
 &
 {\sffamily\textbf L}
 &
Letzte ... (zB letzte Mahlzeit, letzte
Regelblutung, letzter Spitalsaufenthalt)
\\\midrule

{\multirow{8}{2cm}{\includegraphics[width=1cm]{icons/stethoskop.jpg}
BEFUNDE}}
 &
{\sffamily\textbf RR}
 &
auskultatorische und palpatorische Messung
\\
 &
{\sffamily\textbf HF}
 &
Herzfrequenz (Puls), Puls rhythmisch oder arrhytmisch
\\
 &
{\sffamily AF}
 &
Atemfrequenz ausz\"ahlen
\\
 &
{\sffamily $SpO_2$}
 &
Sauerstoffs\"attigung im Blut, Monitoring
\\
 &
{\sffamily  BZ}
 &
Blutzucker {\tiny(Wann war die letzte
Mahlzeit/Insulinspritze? Diabetes bekannt?)}
\\
 &
{\sffamily Temp.}
 &
K\"orpertemperatur/Fieber
\\
 &
 {\sffamily\small Neuro\-status}
 &
Bewusstsein (ev. GCS), Orientierung,
Pupillen, Halbseitenzeichen, Herdblick, Meningismus, retrograde Amnesie
\\
 &
{\sffamily\small Trauma\-check}
 &
abnorme Beweglichkeit, DMS an allen
Extremit\"aten (Nagelbettprobe), Wunden, paradoxe Atmung,
Abwehrspannung des Bauches, ...
\\\midrule
 SANIT\"ATSHILFE

\includegraphics[width=1cm]{icons/san.png}
&
&
Entscheidung Notarzt-Nachforderung

Fortf\"uhrung der Ersten Hilfe

Schienung

Sauerstoffgabe

Reanimationsbereitschaft

Verlaufskontrolle

\ldots
\\\midrule
NOTARZT-ASSISTENZ
\includegraphics[width=1cm]{icons/asb.png}
&
 &
Infusion/peripherven\"osen Zugang vorbereiten

Intubation vorbereiten

Medikamente vorbereiten

EKG anlegen

\ldots
\\\midrule
KRANKENHAUS

\includegraphics[width=1cm]{icons/hospitalStaff.png}
 &
 &
Welches Spital/Bett buche ich ab? (Arbeitsunfall? Welche Abteilung
fahre ich an? Ist der Patient bereits in einem Spital in Behandlung?)

Brauche ich eine Spezialabteilung? (Schockraum, Stroke Unit, PTCA,
Verbrennungsstation)

Arzt\"ubergabe (Anamnese, Befunde, Ma{\ss}nahmen, Hinweise)
\\\bottomrule
\end{tabulary}
\end{table}




\subsection{Ein Beispiel \ldots}
\paragraph[Ein Beispiel: Herr Sepp hat Atemnot]{Ein Beispiel: Herr Sepp
hat Atemnot}


Berufung: Atemnot

\begin{quote}
Sie kommen in eine Wohnung. Die Umgebung erscheint sicher. Die Wohnung
sieht sauber und ordentlich aus. Auf der Couch liegt flach ein
50-j\"ahriger Mann mit weit aufgerissenen Augen und sichtlicher
Atemnot. Er wirkt bla{\ss} und schwei{\ss}ig und fasst ich mit der
rechten Hand an den Brustbereich.

Sie begr\"u{\ss}en den Patienten, stellen sich vor und fragen welche
Beschwerden er hat. Der Patient ringt nach Luft und klagt \"uber
Schmerzen in der Brust. Diese haben vor einer Viertelstunde angefangen.
\end{quote}

Welche Informationen haben Sie bis jetzt gewonnen?
\begin{quote}
Der Patient ist bei Bewusstsein. Er hat Atemnot und einen pl\"otzlich
eingetretenen Thoraxschmerz. Er ist bla{\ss} und schwei{\ss}ig,
vielleicht steckt ein Schockgeschehen dahinter.
\end{quote}
K\"onnen Sie schon etwas tun?
\begin{quote}
Ja.  Bei erhaltenem Bewusstsein, Atemnot und Thoraxschmerz wird der
Patient mit erh\"ohtem Oberk\"orper gelagert. Weiters sprechen die
Symptome f\"ur den Einsatz von Sauerstoff. Es muss noch schnell
gekl\"art werden ob eine Hyperventilationstetanie (nicht
wahrscheinlich) oder COPD-Erkrankung (erfragen) vorliegt.
\end{quote}
Ist der Patient vital bedroht?
\begin{quote}
Ja. Zwar wurden noch keien Vitalwerte erhoben, aber bereits jetzt gilt
der Patient als vital bedroht aufgrund der pl\"otzlich einsetzenden
Thoraxschmerzen und der Atemnot.

\begin{itemize}
    \item Der Oberk\"orper des Patienten wird hochgelagert
    \item Der Notarzt wird nachgefordert (Berufungsgrund: akuter
    Thoraxschmerz mit Atemnot und schlechtem AZ).
    \item Der Patient wird gefragt ob er an COPD leidet, dies wird verneint.
    Eine Sauerstoffmaske mit einem Flow von 10l/min wird angelegt.
    \item Reanimationsbereitschaft wird hergestellt
\end{itemize}
\end{quote}

Was sind die n\"achsten Schritte?
\begin{quote}
Wichtig sind nun die Vitalwerte und das fr\"uhe Beginnen des
Monitorings.Dann muss man genaueres \"uber das Beschwerdebild in
Erfahrung bringen. Der weitere Plan h\"angt davon ab.

Die Vitalwerte werden gemessen und ein Pulsoxy angeh\"angt (RR 100/60,
HF 110, SpO2 94\%). Der Patient gibt an dass er st\"arkste
stechende/brennende Schmerzen hinter dem Brustbein hat die ins Kiefer
ausstrahlen. Ausserdem f\"uhle es sich so an als w\"urde jemand auf dem
Patienten drauf stehen. Die Schmerzen sind nicht atem- oder
bewegungsabh\"angig  und haben vor einer Viertelstunde
{\quotedblbase}wie aus dem Nichts heraus{\textquotedblleft} begonnen.
Ausserdem bekomme er schlecht Luft, sonst h\"atte er keine Beschwerden.
\end{quote}

An welche Differentialdiagnosen denken Sie?

Leitsymptom: Thoraxschmerz 

\begin{tabularx}{\linewidth}[H]{XXXc}
\toprule
\multicolumn{4}{c}{\centering Tabelle: Differentialdiagnosen
Thoraxschmerz unseres Patienten
}\\\midrule
 &
\centering daf\"ur spricht
 &
\centering dagegen spricht
 &
Bewertung
\\\midrule
Herzinfarkt, 

Angina pectoris
 &
Schmerzen hinter dem Brustbein, Ausstrahlung, nicht atemabh\"angig,
Druck-/Engegef\"uhl, Atemnot
 &
 &
sehr wahrscheinlich
\\\midrule
Lungenembolie
 &
Schmerzen hinter dem Brustbein, Atemnot
 &
untypische Ausstrahlung, 
 &
gut m\"oglich
\\\midrule
Pneumothorax
 &
Schmerzen im Brustbereich, Atemnot
 &
pl\"otzliche st\"arkste Schmerzen
 &
eher unwahrscheinlich
\\\midrule
Trauma
 &
Schmerzen
 &
kein Hinweis
 &
unwahrscheinlich
\\\midrule
Erkrankung des St\"utz- und Bewegungsapparats
 &
Schmerzen 
 &
untypische Ausstrahlung, Atemnot, nicht bewegungsabh\"angig
 &
sehr unwahrscheinlich
\\\bottomrule
\end{tabularx}

\gFazit{Die wahrscheinlichsten Diagnosen sind der
Herzinfarkt und die Lungenembolie.}

Fallen Ihnen bei diesen Erkrankungen wichtige weitere Untersuchungen
ein?

Bei beiden Erkrankungen ist es wichtig beide Beine zu begutachten. Im
Falle einer Lungenembolie sucht man nach einer Beinvenenthrombose
(einseitig \"uberw\"armtes, bl\"aulich-r\"otliches geschwollenes Bein).
Bei einer Herzinsuffizienz bzw. Herzinfarkt beurteilt man den
R\"uckstau des Blutes, dh. ob die Beine geschwollen sind.

Ausserdem, hat der Patient so etwas gehabt? Oder sind einschl\"agige
Erkrankungen bekannt?  

\begin{quote}
Sie begutachten beide Beine des Patienten und stellen mittels
Daumendruckprobe fest dass beide Beine leicht geschwollen sind. Auf
Nachfrage gibt der Patient an dass er sonst keine geschwollenen Beine
hat. 
\end{quote}

Was sind die n\"achsten Schritte?

\begin{quote}
Die wichtigsten Untersuchungen wurden gemacht, nun muss man sich um die
Anamnese k\"ummern. 

Symptome: siehe oben.

Krankheiten: Der Patient leidet an nicht-Insulinpflichtigem Diabetes
mellitus und Bluthochdruck. Sonst sind keine Krankheiten bekannt.
Spitalsbriefe gibt es nicht. 

Sie merken sich dass Sie nachher den Blutzucker messen wollen!

Allergien: keine bekannt.

Medikamente: Glucophage 850mg f\"ur den Zucker, Tritace f\"ur den
Blutdruck. Medikamente wurden heute wie verschrieben eingenommen. 

Ereignisse vor Notfalleintritt: Schmerzen begannen pl\"otzlich als der
Patient W\"asche waschen wollte. In den letzten Wochen immer wieder
Stiche in der HErzgegend, hat er aber nicht ernst genommen. 

Letzte(r):

Spitalsaufenthalt: Vor 10 Jahren in der Rudolfstiftung wegen Rotlauf.

Arztbesuch: Vor einer Woche wegen Neuaustellung eines Rezeptes.

Mahlzeit: Mittagessen, Gulasch

Regel: Der Patient ist m\"annlich!
\end{quote}
pro:

Lungenembolie 

Wichtig sind nun die Vitalwerte und das fr\"uhe Beginnen des
Monitorings.Dann muss man genaueres \"uber das Beschwerdebild in
Erfahrung bringen. Der weitere Plan h\"angt davon ab.


\section{Traumatechniken}
Michael Motal

\subsection{Helmabnahme}

immer.

\paragraph*{Material}~

\begin{itemize}
\item 2 Helfer
\item 1 Stifneck f\"ur sp\"ater!
\end{itemize}

\paragraph*{Vorgehen}~

\begin{multicols}{2}
\begin{enumerate}
    \item Ein Helfer fixiert den Helm von hinten und l\"asst nicht mehr los,
    bis der Helm unten ist. Zwischen Helm und Schritt eine Helmh\"ohe
    Abstand lassen, damit der Helm in einer Bewegung abgezogen werden kann
    \item Ein Helfer spricht den Patienten von vorne an, \"offnet Visier und
    Kinnriemen
    \item Der \gE{seitliche} Helfer stabilisiert den Kopf: eine Hand fasst mit
    Daumen und Zeigefinger das Unterkiefer, die andere Hand st\"utzt den
    Kopf von unten. W\"ahrend der Helm abgezogen wird rutscht diese Hand
    etwas nach
    \item Der \gE{hintere} Helfer zieht den Helm in einer S-f\"ormigen Bewegung
    ab: den Helm zuerst kippen bis die Nasenspitze sichtbar ist, dann die
    Unterkante in der gegengleichen Bewegung zum Helfer bewegen
    \item Der \gE{hintere} Helfer zieht den Helm sanft fertig ab und legt ihn
    sicher ab
    \item Der \gE{hintere} Helfer \"ubernimmt den Kopf im doppelten C-Griff:
    Daumen und Zeigefinger um ein Ohr, Kopf leicht auf Zug halten.
    \item Der \gE{seitliche} Helfer legt ein stifneck an (siehe unten)
\end{enumerate}
\end{multicols}

\subsection{Anlegen einer
HWS-Immobilisationsschiene (Stifneck\texttrademark)}

zur Stabilisierung der Halswirbels\"aule bei jedem Verdacht auf eine
Verletzung; auf jeden Fall bei jeder Helmabnahme!

\paragraph*{Material}~

\begin{itemize}
    \item 2 Helfer
    \item 1 stifneck
\end{itemize}

\paragraph*{Vorgehen}~

\begin{multicols}{2}
\begin{enumerate}
    \item Pat. auffordern, den Kopf nicht zu bewegen
    \item nach Helmabnahme: Kopf weiterhin leicht auf Zug halten
    \item der hintere Helfer h\"alt Kopf/HWS in Neutralstellung
    \item der seitliche Helfer misst die Entfernug von Schulter bis
    Unterkiefer in Fingerbreiten aus
    \item stifneck entsprechend einstellen: gemessene Entfernung auftragen:
    Hartplastikrand, der an der Schulter aufliegt unten, Markierung am
    verschiebbaren Teil bzw H\"ohe der Kinnst\"utze oben)
    \item Verschiebemechanismus des stifneck blockieren, ab nun darf kein
    Verstellen mehr m\"oglich  sein
    \item stifneck zuerst am Kinn anlegen/anpassen und festhalten
    \item Klettverschlu{\ss}band so einfalten, dass es nicht am Boden
    scheifen kann (h\"alt sonst nicht mehr)
    \item Nackenteil unter Pat. durchf\"uhren (Vorsicht:
    Klettverschlussband)
    \item fest schlie{\ss}en
    \begin{itemize}
        \item[] Kinnspitze sollte ein kurzes St\"uck \"uber Kante hinausschauen
        (kein Durchrutschen) aber nicht allzu weit (falsch eingestellt)
    \end{itemize}
\end{enumerate}
\end{multicols}


\subsection{Schaufeltrage}

zusammengesetzt aus 2 L\"angsh\"alften; Fu{\ss}teil jeder H\"alfte
l\"angenverstellbar, Kopfteil fixe L\"ange.

Zur beinahe bewegungsfreien Rettung bei Verdacht auf Wirbels\"aulen-,
Becken- oder Oberschenkelverletzungen. Letztendlich soll der Patient
auf der Vakuummatratze gelagert werden. Sollte der Patient gr\"o{\ss}er
als die Trage sein m\"ussen Kopf und Oberk\"orper auf jeden Fall
trotzdem auf der Trage zu liegen kommen. Ausnahme: Beinverletzung, bei
der eine Wirbels\"aulenverletzung ausgeschlossen werden kann.

\paragraph*{Material}~

\begin{itemize}
    \item 2 Helfer
    \item Schaufeltrage
    \item Gurte f\"ur Schaufeltrage
\end{itemize}

\paragraph*{Vorgehen}~

\begin{multicols}{2}
\begin{enumerate}
    \item Schaufeltrage neben Patienten in zusammengebautem Zustand auf die
    richtige L\"ange einstellen. Sollte ein Patient zu gro{\ss} sein
    m\"ussen Kopf und Oberk\"orper auf jeden Fall trotzdem auf der Trage
    liegen! Vorsicht: Markierung der max. L\"ange bei Schraubmodell
    \item \"Offnen der Trage: Schnappverschl\"usse; je nach Methode nur den
    unteren oder beide \"offnen.
    \item Aufschaufeln, 1. Methode: Fu{\ss}seite offen, Pat. wird von oben
    kommend wie auf eine Schere aufgeladen, die Trage unter ihm
    geschlossen. Kontrolle des sicheren Schlu{\ss}es!
    \item Aufschaufeln: 2. Methode: beide Teile getrennt, Pat. leicht an
    Schulter und H\"ufte zur Seite gedreht und jew. Tragenteil
    daruntergeschoben. H\"alften zuerst am Kopf-, dann am Fu{\ss}ende
    schliessen. Kontrolle des sicheren Schlu{\ss}es!
    \item Angurten: 2 Gurte; jeden Gurt zuerst um den Handgriff der einen,
    dann der anderen H\"alfte f\"uhren, dann \"uber dem Patienten
    schlie{\ss}en und festziehen.
    \item Umlagern auf Vakuummatratze: Trage mit vorbeireiteter Matratze
    (s.u.) neben Patienten stellen um diesen einen m\"oglichst geringen Weg
    tragen/bewegen zu m\"ussen
    \item Gurte \"offnen und entfernen.
    \item H\"alften trennen: beide Schnappverschl\"usse \"offnen
    \item H\"alften parallel zum Patienten flach wegziehen. Der Patient soll
    hierbei nicht bewegt werden!
    \item Sollten die Schnappschl\"osser klemmen kann man die Trage wie bei
    der Scherentechnik an nur einer Seite \"offnen und in V-Form (je ein
    Helfer auf einer Seite) entfernen.
\end{enumerate}
\end{multicols}



\subsection{Vakuummatratze}

zur Ruhigstellung des gesamten K\"orpers, bei Verdacht auf m\"ogliche
Wirbels\"aulenverletzungen, Becken- oder Oberschenkelfrakturen (lockere
Indikationsstellung!)

Am Ende soll die Matratze an die K\"orperform der Patienten eng
angepasst sein und somit ein Bewegen verhindern.

\paragraph*{Material}~

\begin{itemize}
    \item 2 Helfer
    \item Fahrtrage
    \item Vakuummatratze
    \item Accuvac
    \item Leintuch
    \item ev. Adapter, je nach Matratzentyp
\end{itemize}


\paragraph*{Vorgehen}~

\begin{multicols}{2}
\begin{enumerate}
    \item Fahrtrage auf tiefste Stufe stellen (ganz unten)
    \item Vorbereiten: Vakuummatratze auf Fahrtrage legen. Die weichere
    Seite ist die Patientenseite, Ventil beim Kopf. Kugeln so verteilen,
    dass der gesamte K\"orper gut sowie die betroffene Stelle besonders gut
    fixiert werden kann.
    \item Vorabsaugen mit Accuvac, bis gut formbar, aber nicht allzu fest.
    \item Vorformen, Ventil verschlie{\ss}en. Wenn vorhanden Leintuch auf
    die Matratze legen.
    \item Angegurteten Patienten mit Schaufeltrage auf die Vakuummatratze
    legen, Schaufeltrage entfernen (siehe
    {\quotedblbase}Schaufeltrage{\textquotedblleft})
    \item Absaugen mit Accuvac, dabei z\"ugig Matratze an den Patienten
    anpassen. Besonderes Augenmerk gilt der Fixierung der verletzten
    Regionen sowie dem Kopf ({\quotedblbase}Ohren{\textquotedblleft}
    formen: Ecken neben dem Kopf hochformen).
    \item Ist kein weiteres Absaugen mehr m\"oglich und die Matratze gut
    angepasst (Matratze muss jetzt bretthart sein) Ventil schlie{\ss}en,
    Accuvac abdrehen.
    \item Pat. an Fahrtrage angurten.
\end{enumerate}
\end{multicols}


\subsection{Vakuum-Beinschiene}

Schienung einer Fu{\ss}-, Knie- oder Unterschenkelverletzung. Achtung:
Oberschenkel: Vakuummatratze

1 Helfer seitlich, 1 Helfer beim Fu{\ss}ende


\paragraph*{Material}~

\begin{itemize}
    \item 2 Helfer
    \item Vakuumbeinschiene
    \item Accuvac
    \item ev. Adapter
\end{itemize}


\paragraph*{Vorgehen}~

\begin{multicols}{2}
\begin{enumerate}
    \item L\"ange anpassen, gegebenenfalls oben umfalten.
    Klettverschlu{\ss}b\"ander herrichten.
    \item Kugeln gut verteilen. Auch am Fu{\ss}ende m\"ussen genug Kugeln
    sein, um dem Fu{\ss} sowohl unter der Sohle als auch am Fu{\ss}r\"ucken
    gut Halt geben zu k\"onnen
    \item Schiene leicht vorabsaugen, damit Kugeln nicht mehr verrutschen
    k\"onnen
    \item Schuh am gesunden Bein ausziehen: Finden der schonensten Methode.
    Ein Helfer h\"alt das Bein auf Zug, der andere \"offnet den Schuh
    (B\"ander ausf\"adeln, zur Not aufschneiden)
    \item vorsichtiges Wiederholen zu zweit am verletztem Bein, Socken
    ausziehen (DMS)
    \item Stiefelgriff: gleich nach Ausziehen von Schuh/Socken nimmt der
    Helfer am Fu{\ss}ende durch das Loch in der Schiene die Ferse des
    verletzten Fu{\ss}es, die andere Hand h\"alt den Fu{\ss}r\"ucken. Das
    Bein wird so auf dem Boden liegend auf Zug gehalten. Der Helfer
    verbleibt so, bis die Schiene fest sitzt!
    \item Der seitliche Helfer richtet die Schiene vorsichtig unter dem
    Patienten aus und fixiert die Schiene mit den Klettverschlussb\"andern.
    Um den Fu{\ss} werden Ohren gebogen und kreuzweise festgeklettet: der
    Fu{\ss} soll von allen Seiten so gut wie m\"oglich umschlossen sein!
    \item Der seitliche Helfer saugt ab und zieht dabei die Klettb\"ander
    immer wieder fest nach. Der andere Helfer h\"alt nach wie vor den
    Fu{\ss} im Stiefelgriff.
    \item Fertig abgesaugt: Ventil schliessen, Accuvac abschalten. Der
    Stiefelgriff kann jetzt losgelassen werden.
\end{enumerate}
\end{multicols}

Die Schiene soll bretthart sein, das Bein leicht auf Zug fixiert. Der
Fu{\ss} soll weder seitlich noch zum K\"orper hin beweglich sein.



\subsection{Sam-Splint}

Unterarmschienung

\paragraph*{Material}~

\begin{itemize}
    \item 2 Helfer
    \item Sam-Splint
    \item Peha-Haft, Mullbinde + Leukoplast oder 2-3 Dreiecktuchkravatten
    \item ein weiteres Dreiecktuch
\end{itemize}


\paragraph*{Vorgehen}~

\begin{enumerate}
    \item Sam-Splint in der H\"alfte falten, an gesundem Arm der L\"ange
    nach anpassen, Knick beim Ellbogen
    \item Anlegen an verletztem Arm, sanft anpassen. Fixierung mit
    Peha-Haft, Mullbinden oder Dreiecktuchkravatten
    \item Fingerspitzen m\"ussen sichtbar sein (DMS), gegebenenfalls Enden
    umfalten.
    \item Lagerung des geschienten Armes im Armtragetuch
\end{enumerate}



\subsection{Druckverband}

zur Stillung einer starken, d.h. potentiell lebensbedrohlichen Blutung

\paragraph*{Material}~

\begin{itemize}
    \item 1-2 Helfer
    \item sterile Wundauflage je nach Wundgr\"o{\ss}e
    \item saugf\"ahiger Druckk\"orper: zB Mullbinde
    \item Dreiecktuchkravatte
\end{itemize}

\paragraph*{Vorgehen}~

\begin{multicols}{2}
\begin{enumerate}
    \item Sofort bei Erkennen einer \gE{starken} Blutung durch
    direkten Druck auf die Wunde und Hochhalten die Blutung verlangsamen (Pat. kann dies je
    nach Zustand selbst tun!)
    \item sterile Wundauflage auf die Wunde legen
    \item Druckk\"orper dem Wundverlauf folgend \"uber die sterile
    Wundauflage legen, weiterhin dr\"ucken.
    \item Druckk\"orper mit einer Dreiecktuchkravatte fest und m\"oglichst
    vollst\"andig umschlie{\ss}en
    \item Dreiecktuchkravatte festziehen und mit einem festen Knoten bei
    ausreichendem Druck direkt \"uber der Wunde fixieren
    \item 2. Knoten zur Sicherung des ersten
\end{enumerate}
\end{multicols}

Sollte die Blutung zu einem sp\"ateren Zeitpunkt wieder
st\"arker werden oder der Druckverband zu locker angelegt
sein wird ein \gEs{weiterer} Druckk\"orper mit einer
weiteren Dreiecktuchkravatte noch fester angelegt. \gEs{Ein
schon angelegten Druckverband wird nicht wieder geöffnet!}






\section{Assistenzt\"atigkeiten}

\begin{tabularx}{\linewidth}[H]{XX}
Teile des Textes und der Handlungsanleitungen enstammen dem Skriptum
{\quotedblbase}\"Arztliche Grundfertigkeiten{\textquotedblleft} der 
Besonderen Einrichtung zur medizinischen Aus- und Weiterbildung --
Medizinische Universit\"at Wien, Abt. Methodik und Entwicklung, welches
in Zusammenarbeit  mit dem Klinischem Institut f\"ur Hygiene und
Medizinische Mikrobiologie entstanden ist. 

Wir danken Hrn. ao.Univ.Prof. Dr. Michael Schmidts und
seinem Team f\"ur die Unterst\"utzung!
 &
Powered by
Image
% Unhandled or unsupported graphics:
%\includegraphics[width=4.551cm]{icons/muw.png}
\\
\end{tabularx}


% ========Subsection eingefügt vom Emhofer 1.5.2009===============================

\subsection{\"Ubergabegespr\"ach an Notarzt bzw. Spitalspersonal}  

Die \"Ubergabe des Patienten in den Verantwortungsbereich eines Notarztes oder des Spitalspersonals hat geordnet und gewissenhaft zu erfolgen, damit keine wichitge Information verloren geht. Diese Schnittstelle zwischen Sanit\"ater und Arzt bzw. Krankenschwester kann auf den weiteren Behandlungsverlauf einen großen Einfluss nehmen und darf daher in seiner Bedeutung nicht untersch\"atzt werden! Es empfiehlt sich -- so Zeit dafür ist -- sich bereits vor Eintreffen des Notarztes oder vor Ankunft im Krankenhaus zu \"uberlegen, welche Informationen wesentlich sind und daher auf keinen Fall vergessen werden d\"urfen.

Das \"Ubergabegespr\"ach muss mindestens folgende Informationen beinhalten:
\begin{itemize}
    \item Name des Patienten
    \item Alter
    \item Zusammenfassung der wesentlichen Inhalte der Anamnese
    \item Beschreibung der Symptomatik und deren Verlauf
    \item bisher erhobene Befunde (Blutdruck, Puls, Sauerstoffs\"attigung usw.)
    \item bisher durchgef\"uhrte Maßnahmen
    \item sonstige, wichtige Hinweise (z.B. Infektionsgefahren, Bluterkrankheit etc.)
\end{itemize}

%=========Ende Abschnitt Emhofer==========================================================


\subsection{Intubation: Vorbereitung und Assistenz}

\paragraph*{Material}~

\begin{multicols}{2}
\begin{compactenum}
    \item Beatmungsbeutel inkl. passender Beatmungsmaske, Bakterienfilter und
    Sauerstoffreservoir
    \item $O_2$-Beriselungseinheit und $O_2$-Line
    \item Evt. Beatmungsgerät
    \item Absaugeeinheit: \gCa{Keine Intubation ohne Absaugung!}
    \item Laryngoskop, bestehnd aus
    \begin{compactenum}
        \item Spatel: gebogen oder gerade; Größen 0--4
        \item Griff
    \end{compactenum}
    \item (Endotracheal-)Tubus; verschiedene Größen
    \item Führungsdraht (\gT{Mandrin})
    \item Silikonspray
    \item Evt. Magillzange
    \item Mit Luft gefüllte Spritze zum Cuffen (10\,ml)
    \item Beißschutz: Beißkeil oder Guedel-Tubus
    \item Fixationsmaterial (Mullbinde)
    \item Stethoskop
\end{compactenum}
\end{multicols}


\paragraph*{Vorbereitung}~

\begin{multicols}{2}


\begin{enumerate}
    \item Gewünschte Spatelgröße erfragen
    \item Gewünschte Tubusgröße erfragen
    \item Laryngoskop (Spatel und Griff) zusammenbauen
    \item Tubusverpackung aufschälen, Tubus aber in der Verpackung belassen
    \item Führungsdraht (Mandrin) einführen. Die Spitze muß bis zum Tubusende
    vorgeschoben vorgeschoben werden, darf aber nicht herausragen.
    \item Cuff mit Silikonspray einführen
    \item Stethoskop: Bei Doppelkopfstethoskopen Membran-Seite einstellen.
    \item Absaugung bereit machen:
    \begin{enumerate}
        \item Passenden Absaugkatheter wählen
        \item Verpackung teilweise aufschälen um Anschluß freizulegen, Katheter
        in der Verpackung belassen.
        \item Katheter anschließen.
        \item Funktionskontrolle: Gerät einschalten. Saugstärke?
        Batteriewarnung?
\item Absaugeinheit muß in Griffweite der Assistenten positioniert werden! 
    \end{enumerate}
    \item Beatmungsbeutel zusammenbauen und an $O_2$-Beriselungseinheit
    anschließen
    \item Restliches Material in Nierentasse bereitlegen
\end{enumerate}
\end{multicols}

\paragraph*{Assistenz}~
\begin{multicols}{2}


\begin{enumerate}
    \item Überprüfen ob alle Materialen bereit sind.
    \item Sauerstoffberiselung einschalten und auf 15\gLMin einstellen,
    Reservoir des Beatmungsbeutels füllen. 
    \item Durchführenden melden daß alle Materialen bereit sind.
    \item \gZ{Evt. werden jetzt von einem Arzt Medikamente verabreicht die
    die Intubation ermöglichen. Der Durchführende wird nun einige Zeit den
    Patienten mittels Beutel-/Masken-Beatmung beatmen (sofern nicht bereits
    geschehen).}
    \item Laryngoskop in die \gEs{linke Hand} zureichen\footnote{Es ist egal ob
    der Durchführende Rechts- oder Linkshänder ist, das Laryngoskop wird
    \gE{immer in der linken Hand} gehalten.}
    \item Auf Anweisung \emph{"`Kehlkopfdruck"', "`Krickoiddruck"'} auf den
    Kehlkopf drücken.
    \item Tubus in die Rechte Hand zureichen.
    \item Auf Anweisung saugen.
    \item Beatmungsmaske vom Beatmunsbeutel trennen.
    \item \gZ{Durchführender intubiert und entfernt Führungsdraht (Mandrin)}
    \item Mit luftgefüllter Spritze Cuff aufblasen (\gE{"`cuffen"'})
    \item Beatmungsbeutel an Tubus anschließen
    \item Notarzt Stethoskop in die Ohren geben. 
    \item Es erfolgen 3 Atemhübe wobei das Sthetoskop über Magen, linkem
    Lungenflügel und rechtem Lungenflügel positioniert wird. Es wird dabei
    überprüft ob der Tubus richtig liegt. 
    \item Der Druchführende übernimmt die Beatmung. 
    \item Einführen des Beißkeils oder des Guedel-Tubus
    \item Fixierung des Beißschutzes und des Tubus mittels Mullbinde
    \item \gZ{Ggfs. BEatmungsgerät einstellen und
    anschließen}
\end{enumerate}
\end{multicols}
\begin{enumerate}
    \item 
\end{enumerate}

\gMarried
{\gAuthNotes{Text fehlt}}
{
\includegraphics[width=\linewidth]{./images/intubation_zubehoer_op.jpg}
\captionof{figure}{Intubationszubehör. Angerichtet auf einem Anästhesiegerät.
\gSourcePicture{Gabriel}
\gAuthNotesPicLegalOk{intubation\,zubehoer\,op.jpg}{cc-by}}
}

\subsection{Aufziehen eines Medikamentes aus einer Ampulle}

Ampulle mit rotem Punkt: Sollbruchstelle, Daumen muss beim aufbrechen
auf dem roten Punkt sein. 

Ampulle mit {\quotedblbase}Halsring{\textquotedblleft}: Ampulle kann
nach alle Richtungen aufgebrochen werden. 

\marginpar{\gH{Warum mach{\textquotesingle} ich mir wegen
der Verletzungsgefahr Sorgen? Ist ja nur ein kleiner
Schnitt, ist ja meine Sache {\dots}}

Nein! Unter Normalbedingungen, wenn ich in aller Ruhe mein Medi
aufziehen kann ist es tats\"achlich mein eigenes Bier. Was aber wenn
ich mich bei einer Reanimation an der Ampulle schneide? Die anderen
Kollegen sind besch\"aftigt (Dr\"ucken, blasen, Defi {\dots}) und dann
steh ich da mit meinem kleinen Schnitt und blute alles an. Und der
Notarzt braucht dringend das Adrenalin {\dots} Was jetzt?

Inakzeptabel!}

\paragraph*{Material}~ 

\begin{itemize}
    \item Spritze
    \item Kan\"ule(n)
    \item Ampulle mit Medikament
    \item Trockene Tupfer
\end{itemize}

\paragraph*{Aufziehen des Medikamentes}

\begin{multicols}{2}

\begin{enumerate}
    \item Kontrolle Ablaufdatum
    \item Kontrolle Inhalt: Verf\"arbung, Ausflockung?
    \item \"Offnen der Ampulle
    \begin{itemize}
        \item  Glasampulle
        \begin{itemize}
            \item Ampullenspitze beklopfen, darauf achten dass sich keine
            Fl\"ussigkeit im Apullenkopf mehr befindet
            \item \"Offnen der Medikamenten-Ampulle mit Tupfer, roter Punkt weist
            zum K\"orper
\marginpar{\includegraphics[width=\linewidth]{./images/ampulle_beklopfen.png}\captionof{figure}{Ampullenkopf
beklopfen. \emph{Bild: BEMAW}}}
        \end{itemize}
        \item Plastikampulle,
        {\quotedblbase}Dreh-und-Drink{\textquotedblleft}-Verschlu{\ss}
        \begin{itemize}
            \item Den Verschlussknopf ohne Kontamination des Ampullenhalses
            (Sollbruchstelle) abdrehen
        \end{itemize}
    \end{itemize}
    \item Ampulle abstellen
    \item Verpackung der Spritze aufsch\"alen, Spritze in Verpackung lassen
    \item Verpackung der Nadel aufsch\"alen, Spritze auf Konus aufstecken
    \item Entfernen der Kan\"uleabdeckung (nicht abdrehen, sondern
    abziehen!)
    \item Einf\"uhren der Nadel in die Ampulle
    \begin{itemize}
        \item Beim Einf\"uhren der Kan\"ule in die Ampulle den Ampullenrand
        nicht ber\"uhren
    \end{itemize}
    \item Aufziehen des Medikamentes
    \item Vorsichtig entl\"uften
    \item Variante {\quotedblbase}Verabreichung \"uber
    Venflon{\textquotedblleft}
    \begin{itemize}
        \item Nadel h\"andisch abziehen und in den Kontamed verwerfen.
        \item Arzt das Medikament \"ubergeben, Ampulle zeigen.
        \item Wird das Medikamen nicht sofort verwendet spritze mit roten
        Stoppel verschliessen und Ampulle ankleben.
    \end{itemize}
    \item Variante {\quotedblbase}Zuspritzen in eine
    Infusionsflasche{\textquotedblleft}
    \begin{itemize}
        \item Gummimembran desinfizieren
        \item Kan\"ule nicht bis zum Schaft einstechen
        \item Medikament zuspritzen
        \item Beschriften der Flasche
    \end{itemize}
    \item Variante {\quotedblbase}direktes Spritzen{\textquotedblleft}
    \begin{itemize}
        \item Arzt nach Nadelgr\"o{\ss}e fragen
        \item Aufziehnadel h\"andisch entfernen
        \item Spritznadel aufstecken
        \item Arzt Spritze mit Verschlusskappe reichen, Ampulle zeigen.
    \end{itemize}
\end{enumerate}
\end{multicols}

\subsection{Vorbereitung einer Infusion}

\paragraph*{Material}

\begin{itemize}
    \item Infusionsflasche
    \item Infusionsbesteck
    \item Aufh\"angevorrichtung
    \item alkoholhaltiges Desinfektionsmittel
    \item einige saubere Tupfer
    \item Permanentmarker
\end{itemize}

\paragraph*{Handlungsschritte}~
\marginpar{\includegraphics[width=\linewidth]{./images/infusion_tropfkammer.jpg}}
\begin{multicols}{2}
\begin{enumerate}
    \item H\"andedesinfektion
    \item Aufh\"angevorrichtung auf Flasche aufsetzen
    \item Er\"offnen der Flaschenabdeckung (Membranschutz)
    \item Wischdesinfektion der Gummimembran
    \begin{itemize}
        \item Alkoholgetr\"ankter Tupfer kann w\"ahrend weiterer Vorbereitung
        auf Membran belassen werden
    \end{itemize}
    \item Montage des Infusionsbestecks
    \begin{itemize}
        \item \gEs{Radklemme schlie{\ss}en}
        \item Einstich des Tropfkammerdorns in Infusionsflasche mit
        geschlossenem L\"uftungsventil und geschlossener Radklemme
        \item {\quotedblbase}No touch technique{\textquotedblleft} -- Dorn der
        Tropfkammer und Gummimembran der Infusionsflasche d\"urfen nicht
        ber\"uhrt werden
    \end{itemize}
    \item Infusionsschlauch f\"ullen
    \begin{itemize}
        \item Flasche wenden und hochhalten
        \item Tropfkammerluftventil \"offnen
        \item Tropfkammer durch Zusammendr\"ucken bis zur H\"alfte f\"ullen
        \item Konusabdeckung des Infusionsschlauchs am anderen Ende muss bei der
        Bel\"uftung nicht entfernt werden
        \item Infusionsflasche und Ende des Infusionsschlauchs mit nicht
        dominanter Hand halten  \gZ{(Wenn sich das
        Infusionsschlauchende auf gleicher H\"ohe wie der Spiegel in der Tropfkammer befindet,
        entl\"uftet sich der Schlauch problemlos --  Prinzip der
        kommunizierenden Gef\"a{\ss}e, an kommunizierenden Gef\"a{\ss}en stellt
        sich immer dieselbe Fl\"ussigkeitsspiegelh\"ohe
        ein.)}
        \item Radklemme \"offnen
        \item Infusionsschlauch {\quotedblbase}ohne
        vergie{\ss}en{\textquotedblleft} entl\"uften
        \gZ{(Fl\"ussigkeit bis ca. 5cm vor Konnektor rinnen
        lassen, endg\"ultig entl\"uften bei Patienten vor
        Verabreichung)}
        \item Radklemme schlie{\ss}en
    \end{itemize}
\end{enumerate}
\end{multicols}


\subsection{Venenverweilkan\"ule}
({\quotedblbase}Venflon{\textquotedblleft}): Vorbereitung und Assistenz 

\begin{figure}
\begin{center}
\includegraphics[width=13.467cm]{./images/venflon.png}
\end{center}
\caption{Eine periphere Venenverweilkanüle (Venflon\textregistered), in ihre
Einzelbestandteile zerlegt. \gSourcePicture{BEMAW}}
\end{figure}

\begin{table}
\begin{tabular}{m{1.5619999cm}m{1.7229999cm}m{1.7229999cm}m{2.0509999cm}m{1.7839999cm}m{2.102cm}}
\toprule
\multicolumn{3}{m{5.408cm}}{Nur zur Illustration!} 
&
 Durchfluss [ml/min]
 &
 Farbcode
 &
\\
 {\O} mm & G & Lmm  & & &
\\\midrule
0,5 x 0,8 & 22 & 25 & 25 & blau & Sehr d\"unn
\\
0,7 x 1,0 & 20 & 32 & 55 & rosa & d\"unn
\\
0,9 x 1,2 & 18 & 40 & 90 & gr\"un & normal
\\
1,1 x 1,4 & 17 & 42 & 135 & wei{\ss} & dick
\\
1,3 x 1,7 & 16 & 42 & 170 & grau & Sehr dick
\\
1,6 x 2,1 & 14 & 42 & 265 & orange & Pferd
\\\bottomrule
\end{tabular}
\caption{Verschiedene Größen von peripheren
Venenverweilkanülen. \gSource{BEMAW (modifiziert)}}
\end{table}


\paragraph*{Material}~

\begin{multicols}{2}
\begin{itemize}
    \item Nierentasse
    \item Staubinde
    \item Venflon:
    \begin{itemize}
        \item Die Gr\"o{\ss}e (= Farbe) ist vom Arzt vorher zu erfragen!)
        \item Die Fl\"ugel werden hinunter geklappt
        \item Die farbige Lasche des Zuspritzventils in rechten Winkel zur
        Venflonachse stellen.
    \end{itemize}
    \item Einmalverschlusskappe ({\quotedblbase}roter
    Stoppel{\textquotedblleft})
    \item Fixierungspflaster
    \item einige saubere trockene Tupfer
    \item Alko-Tupfer
    \item (Einwegspritze mit 5ml Sp\"ulfl\"ussigkeit)
    \item (Faserschreiber)
    \item Kontamed-Box
\end{itemize}
\end{multicols}

\paragraph*{Assistenz}~

\begin{itemize}
    \item Stauschlauch zureichen
    \item ge\"offnete Verpackung des Alko-Tupfer hinreichen sodass dieser
    entnommen werden kann
    \item Kontamed in Reichweite des  Arztes\,/\,NKV
    bereitstellen sodass er sp\"ater die Nadel abwerfen kann.
    \item Venflon zureichen
\end{itemize}
\begin{itemize}
    \item Venflon an der Schutzkappe halten sodass der
    Arzt\,/\,NKV den Venflon abziehen kann.
    
\end{itemize}
\begin{itemize}
    \item Trockene Tupfer zureichen.
    \item Je nach Situation Infusionsschlauch, roten Stoppel oder gar nichts
    zureichen.
    \item Fixierpflaster zureichen
\end{itemize}
